\section{Additional benchmark qualification}
\label{sec:incremental-recruitment}
An independently pinned SWE-bench Verified metadata projection~\citep{swebench2024,swebenchverified2024} supplies 500
candidate repairs from twelve repository names. Ten of eleven prospectively
selected two-issue pairs have forward base-commit ancestry; one has reverse
ancestry in PR-date order. A separate current GitHub identity audit checks all
192 candidate/original repository pairs, finds no shared repository or declared
fork identities, and flags two shared-owner groups. This does not establish
historical code independence, portfolio separation or absence from model
training. A later audit reports test-specification problems in a selected hard
subset and evidence of solution exposure in tested frontier models~\citep{swebenchverifiedaudit2026}.
Its selected-subset rate is not a whole-dataset error estimate, and it does not
measure contamination in our provider profile. Public examples exposed during
this literature review are separately recorded before any prospective split.
No new confirmatory family or complete natural trajectory is enrolled.

The first Requests node has a separately named ARM environment preserving its
public Python version and thirteen listed Python package versions. The unchanged
native SWE-bench 5.0.2 parser/grader observes all 187 designated checks in both
controls: unfixed source passes 185 and fails the two required failing checks;
reference source passes all 187 and is the only resolved control. Both collected
outputs also contain eight missing-fixture setup errors outside the designated
set. The original script lacks a test-exit marker, so a zero container exit
cannot establish acceptance. This is designated-check qualification, not an
error-free full-suite result or equivalence to the published x86 environment.
Protected reference/test content remains in curator custody; model-facing
selection reads public metadata only. Earlier import and decoding failures are
retained. Infrastructure qualification contributes no provider-token or repair
success observation for either treatment.

The second Requests node subsequently qualifies in a new named ARM profile
using its public development dependencies and Docker's default bridge network.
Both controls complete all 131 designated checks: unfixed source passes 130
and the reference passes all 131. The full parsed outputs have 288 versus 289
passes, respectively, with one expected failure each and no setup errors.
The original offline port, a dependency-build failure and a plugin-startup
failure remain retained; the successful profile is a development correction.
This does not make the original failed port successful. A later candidate
grading service version fails its four-control qualification: three controls
behave as expected, but this node's reference is not accepted despite all
designated checks completing. Four existing public network-timeout checks
instead encounter a network-unreachable condition; the earlier standalone
run used the same image, package versions and evaluation bytes. The unique
cause and a reliable remedy are not established by this discrepancy.
Eight service rejection controls pass. This retained discrepancy illustrates
why standalone benchmark qualification alone does not qualify a new endpoint.
A subsequently frozen service version passes one attempt at each of its four
native controls and all ten rejection controls. It preserves exact objective,
profile, cumulative patch and resulting source-tree bindings. Public network
readiness probes bracket the first node's evaluation, with no automatic retries;
an incomplete grade remains unavailable for repair decisions. All eight
qualification probes pass. The 94-check audit retains 1,175 artifacts; the four
control containers take 87.429894 seconds in total. This qualifies the named
service and its observed controls, without proving the unique cause or remedy
of the previous network failure or stability throughout future test runs.
Candidate and benchmark tests still share the original test process; adversarial
oracle integrity is not established. The second node retains its eight
outside-set fixture errors.

An additive public-only cumulative-source control changes two production files
and adds one Git-ignored Python file. The service inventories all 98 resulting
files, matches the independently expected tree, and returns the unchanged
unfixed grade of 185/187 in 6.048 seconds. This checks complete source-carry
transport into grading, not successful repair or an accepted multi-issue
trajectory. Earlier four-control cohorts remain unchanged.

A further pinned metadata-only projection yields 21,336 SWE-rebench
rows~\citep{swerebenchmetadata2026} from 3,468 repository names and 731
SWE-bench Pro rows~\citep{swebenchprometadata2026} from eleven names.
The 22,067-row sum is not deduplicated across datasets; protected issue,
reference and test fields are not decoded in this projection. A prospectively
seeded screen uses 2025 metadata dates, available image fields, distinct base
commits and exclusion of prior repository/owner names. It yields 53 candidate
pairs. All first sixteen pairs have actual forward Git ancestry, and 448
comparisons against the 28 prior current repository identities find no
repository, declared-fork or owner match. Each candidate has a distinct current
repository and owner identity. These checks do not prove historical code
independence, portfolio separation or training novelty. Image-field eligibility
is a convenience filter; all 32 selected image names specify x86\_64, whereas
the evaluation host is ARM. Source/environment/scorer qualification and a
prospective confirmatory split remain open; zero new families are enrolled.

For a later development experiment on the existing pyparsing snapshot, the
current public issue report supplies a concrete objective. Its nine examples
reproduce the reported recursive-whitespace defect on unchanged public source:
eight conform to the reported expectation and one fails. The report is frozen
for that new experiment; its exact bytes at the original base date remain
unverified. A reproduced public example is not complete hidden acceptance.

\section{Repository preparation and source custody}
\label{sec:incremental-context-preparation}
\paragraph{Repository context custody and preparation cost.}
A separate native producer/consumer qualification retains all 302 regular files
across the XML and pyparsing development snapshots, including unparsed source.
Its bounded semantic scope analyzes one of six and five of 154 Python files,
respectively; four dependency capsules grant no source omission. Four
cold/incremental root-and-block comparisons agree, four exact source expansions
execute, and 22 negative controls reject. An added caller checks required-source
CIDs and current bytes because native freshness omits that mapping; native
``sufficient'' flags are not behavioral authority. Full pyparsing scanning
produces 3,957 symbols and 55,859 edges, but the unmodified producer process does
not finish within its 600-second bound. Hoisting a repeated immutable
node-CID index preserves the full XML state/root and reduces its preparation
from 64.838 to 26.296 seconds in one observation per profile. Full pyparsing
still reaches the 600-second bound after that change. Source retention and
bounded integration therefore do not establish affordable whole-repository
preparation, context compression or repair efficacy. Three later isolated
profiles add a fixed-profile CID encoder, cached syntax validation and an exact
literal-union predicate. Across all five profiles XML state/root identity is
unchanged and every native consumer passes; each full pyparsing producer still
hits the process deadline. In the last profile the producer writes its result
at 597.222 seconds before the 600.202-second termination; the native consumer
does not run. An earlier stage diagnostic had also assembled and verified a
full bundle before timing out in source-binding reads. Missing final output
therefore does not prove that bundle construction itself failed.

The literal-union change preserves all 55,859 actual link predicates against
3,957 node CIDs, an embedded-CID rejection, 1,235 constructed comparisons and
64 native regressions. The independently validated full bundle has 67,872
blocks and agrees with the last producer's state/root. Its 72.358 versus 2.834
seconds measure only one diagnostic predicate pass per implementation on a
shared host; they are not a whole-pipeline speedup estimate.

A generic current-repository service separately checks exact source inventory,
file modes and objective/request identities. Twelve controls pass across three
XML preparations: a source change requires new state/root/pack identities,
whereas an objective-only change retains state/root and creates a new pack.
Both public Requests nodes then complete actual native producer/consumer
roundtrips, retaining 97 files each and scanning all 35 declared Python files.
Eleven negative controls pass per node; neither node permits any raw-source
omission. These are context and custody checks, with zero provider calls and
no repair-quality or compression comparison.
Four subsequent Requests API expansions return 79,942 exact raw bytes across
two real read processes, preserving source, request, freeze and objective
identities. They add no omissions or model calls.

A sixth isolated profile memoizes identity only within an exact immutable
native state instance under frozen codec/runtime bindings. Both full producer
processes and their separate native consumers now complete. XML preparation
takes 9.103 seconds; pyparsing takes 566.254 seconds, with
553.516 CPU seconds. The latter producer process takes 569.609
seconds and its consumer 1.740 seconds. All 302 regular files and
11,781,935 bytes are materialized; pyparsing's additional opaque symlink is
retained without following it. Both sources' previous available state/root
identities agree, including the independently checked 67,872-block pyparsing
bundle. All 22 binding/tampering controls pass and all 7,255 frozen runtime
source files remain unchanged. Zero source omissions and zero signed receipts
are retained, along with the unresolved semantic frontier.

The profile passes 159 native regressions. Separate full XML unchanged,
comment-change cold and comment-change incremental states preserve serialized
identities and all 1,992 reachable blocks across the parent and memoized
profiles. Each profile performs twenty actual source materializations and
rejects twenty corrupt-source reads. This qualifies immutable-object identity
caching, not execution/proof reuse or a defense against reflective mutation of
frozen objects. Full pyparsing incremental mutation equality remains separate
work. The successful full run establishes completion within the bound, not a
replicated speedup or token benefit; the five previous deadlines remain results.


\section{Combined controller, publication and recovery}
\label{sec:incremental-controller-publication}
\paragraph{First combined controller qualification.}
A new caller joins actual native context reads, source-bound semantic analysis,
symbolic routing, compiler-mediated world memory, isolated behavioral scoring
and a durable effect journal. On the four existing constructed routing tasks,
two supported initial repairs pass and two remain residual-pending with the
provider disabled. Authored follow-up events consume accepted source, recheck
a comment change and reject a stale local parent. Ordinary no-op detection
receives the same current behavioral checks as the memory path.
Three actual process interruptions exercise validated continuation, durable
receipt reconciliation and an executed effect whose completion remains unknown;
the unknown effect is held without redispatch. Across thirteen events there are
nine local acceptances, two pending residuals, one stale parent and one unknown
effect. These event counts are not thirteen independent repair tasks. Thirty
scorer attempts yield 29 durable terminal receipts. Local ledger CAS is not
native Kit publication; sealing, lifecycle reuse and certified procedures are
not yet combined here. With zero provider calls this qualification establishes
no integrated token-savings effect.


\paragraph{Subsequent native publication profile.}
The separately frozen eight-event sequence yields six acceptances, one stale
branch and one pending residual. A full checkpoint at generation 0 is followed
by a delta at generation 1. A real service exits with code 94 after native
seal persistence and before CAS. A fresh process loads the pending seal,
rechecks the entire parent chain and exact local evidence, replays visible
and audit behavior, then completes native CAS/WAL through an explicitly added
development adapter, reaching generation 2. Another fresh event consumes the
native stored accepted source and publishes generation 3. A branch based on
generation 0 reaches native publication and loses with a generation mismatch.

\begin{table}[htbp]
\centering
\caption{Actual work in the eight-event native publication qualification.
Event outcomes are dependent continuation observations.}
\begin{tabular}{lr}
\toprule
Operation & Count \\
\midrule
Controller invocations / terminal effects & 9 / 68 \\
Native context reads / source expansions & 4 / 8 \\
Input scans / accepted-candidate scans & 8 / 3 \\
Symbolic attempts / produced candidates & 4 / 3 \\
Native routing calls & 8 \\
World-memory service calls / compiled actions & 11 / 2 \\
Publication attempts / native reads / recovery & 7 / 4 / 1 \\
Controller / additional admission scorer processes & 18 / 24 \\
Provider requests / provider tokens & 0 / 0 \\
\bottomrule
\end{tabular}
\end{table}

Invocation wall time is 60.906438 seconds; the additional 24 behavioral
replays consume 13.759197 seconds within that total, not as a separate additive
cost. Context enrollment, inventory freezing, implementation and read-only
audit are outside this runtime denominator. Crash-time work is retained through
a durable receipt, and recovery does not double-count its aliased completion.
Both ordinary and memory paths run the same cheap current-source check and
accept no-op revisits without symbolic or model repair. The compiled memory
action dispatches an actual check; no reduction in checks is demonstrated.

The complete required leaf set is externally frozen: accepted source,
semantic preparation, effect inventory, visible behavior and audit behavior.
Every changed candidate receives a new native scan. All five leaves are
replaced per delta. Native integrity verification checks content and digests;
unsigned local receipts additionally depend on trusted controller/scorer
execution, a durable effect ledger and an evaluation-owned CID index. This
is neither a theorem about arbitrary Python nor production trust admission.
Stored parent chains are recursively checked and the requested source is
behaviorally replayed on read. Broader crash phases, concurrent publication,
provider effects, real repository edits and lifecycle reuse remain open.

The read-only audit passes 6,424 checks. The archive retains 504 source and
metadata files, including 483 native source files across the loaded closure
and preceding context enrollment. The first profile's omitted imported wrapper
and the first audit's pointer-envelope parser failure are retained. Correcting
the audit did not rerun behavioral outcomes; the archived previous controller
still passes its original 9,188 checks.

\paragraph{Full-repository service fixture and measured cost.}
A separate curator-held Requests 6028 fixture begins with its reference
candidate and then adds a neutral comment to the same accepted production
source. Both states contain all 97 source files. It executes two full native
context preparations, four world-memory service calls and six original native
grader invocations: two initial checks and four fresh publication, interruption,
recovery and read replays. All six complete the 187 designated checks.
Reference and hidden test payloads remain in curator custody. This fixture
does not execute the Requests 5414-to-6028 upstream transition or any model
repair, and it does not establish adversarial oracle integrity.

The first full checkpoint reaches generation 0. The second process constructs
a delta and actually exits with code 94 after seal persistence but before CAS.
A fresh process recovers generation 1; another fresh process reads and restores
the exact second accepted source. Wrong-tree, wrong-receipt, source-drift and
incomplete-effect controls all reject before grader replay. All four required
leaves---accepted source, semantic preparation, observed non-sealing effects
and native behavior---are replaced on transition. The effect leaf excludes
all sealing operations, including prior reads; the outer workflow ledger must
still count their actual work. Parent/source admission and fresh grading are
additional checks. Unsigned local receipts and native integrity mode establish
neither a theorem of execution nor verification reuse.

\begin{table}[htbp]
\centering
\caption{Measured disjoint sequential intervals for the Requests 6028 curator
fixture. Grader processes are nested in these phases and are not added again.}
\begin{tabular}{lrr}
\toprule
Phase & Count & Wall seconds \\
\midrule
Full native context preparation & 2 & 318.283 \\
Initial accepted-candidate checks & 2 & 16.047 \\
World-memory service calls & 4 & 5.858 \\
Publication/recovery/read and rejection controls & 8 & 37.037 \\
\midrule
Sum of measured intervals & & 377.226 \\
\bottomrule
\end{tabular}
\end{table}

The six grader containers consume 31.625 seconds within the measured
intervals. Their four replay service wrappers are also nested, not independent
extra costs. The frozen-start-to-result artifact span is 377.640 seconds;
this is not a separately instrumented whole-process timer and excludes earlier
imports and profile creation. Later code archival and read-only audit take
0.275 seconds together, outside the fixture intervals. There is one
shared-host observation, no parallel comparison and no provider calls.
The 275-check audit retains 33 public code snapshots and hashes 1,136
private artifacts. This closes the stated service controls while leaving
multi-issue model efficacy, token savings and broader crash coverage open.

\section{Native proof profiles and exact statement checks}
\label{sec:incremental-native-proof}
\paragraph{A real native proof relation.}
The copied native backend was built offline with locked dependencies because
the bundled binary did not implement the source's v3 circuit. Isolated
development keys support a four-leaf Event-DAG Merkle/count relation. Three
actual proofs for counts 1, 2 and 4 verify; eight witness, public-input,
proof and key mutations reject. Primary setup takes 7.810 seconds,
proof generation 12.048--12.474 seconds, and verification
3.053--5.147 milliseconds; serialized proof JSON is 1,484 bytes.
The raw verifier accepts a different outer root that maps to the same field
element. The evaluation caller's exact-root/count and key/binary checks reject
that alternate statement before dispatch. This boundary is part of the result:
the native relation alone does not bind arbitrary outer metadata. Development
setup is not production trust admission, and a Merkle relation is not proof
that arbitrary tests executed. The v5 pytest-attestation circuit remains
unsupported in this backend.

\paragraph{The original bounded Horn profile.}
The source-built v2 backend also produces and verifies three honest implication
chains of lengths 1, 2 and 4. Setup takes 1.101 seconds, proving
2.238--2.390 seconds, and verification
3.187--5.267 milliseconds; each proof JSON is 1,700 bytes.
The honest prover rejects the tested incomplete or inconsistent inputs.
When the scalar public inputs remain unchanged, the raw verifier accepts three
changed outer metadata fields; exact statement and source/key checks at the
evaluation caller reject them before dispatch. Honest examples and wrapper
checks do not establish that every satisfying circuit witness corresponds to
the advertised textual axioms. That stronger assurance remains unconfirmed.

